\documentclass{ximera}
\begin{document}

    \title{Lecture 19}
    
    \begin{abstract} 
        - Finish Cholesky factorization \\
        - Finish options for Least Squares solution methods
    \end{abstract}
    \maketitle

\textit{This lecture is a combination of Lecture 23 and the part of Lecture 12 we did not finish}

Cholesky factorization is a clever way to do Gaussian elimination on a nicely structured matrix.  Based on what we did last time, we can write (almost) any HPD matrix $A = \begin{bmatrix}1 & w_1^* \\ w_1 & K_1\end{bmatrix}$, where $K_1$ is HPD.  Then, per our work in class last week,
\begin{align*}
    A = \begin{bmatrix}
        1 & 0 \\ w & I
    \end{bmatrix}
    \begin{bmatrix}
        1 & 0 \\ 0 & K-ww^*
    \end{bmatrix}
    \begin{bmatrix}
        1 & w^* \\
        0 & I
    \end{bmatrix}
    := LDL^*
\end{align*}

Here, the first matrix on the left is the $L$ in the $LU$ decomposition of $A$, whereas $DL^* = U$, due to the symmetric structure of the matrix.  Note the suggestion that $D$ should *eventually* be diagonal, while right now it is only block diagonal.  

What about a general HPD martrix $A$?  We can write this as $A = \begin{bmatrix}a_{11} & w_1^* \\ w_1 & K_1\end{bmatrix}$, where $K_1$ is HPD.  Indeed, $a_{11} > 0$ since $e_1^TAe_1 = a_{11} > 0$ by positive definiteness.  This argument can be repeated for all diagonal entires of an HPD matrix.  The corresponding base induction step is therefore
\begin{align*}
    A &= \begin{bmatrix}
        a_{11} & w_1^* \\ w_1 & K_1
    \end{bmatrix} 
    \\
    &=
    \begin{bmatrix}
        \sqrt{a_{11}}& 0 \\ \frac{1}{\sqrt{a_{11}}}w_1 & I
    \end{bmatrix}
    \begin{bmatrix}
        1 & 0 \\ 0 & \frac{1}{a_{11}}(K-ww^*)
    \end{bmatrix}
    \begin{bmatrix}
        \sqrt{a_{11}}& \frac{1}{\sqrt{a_{11}}}w_1^*  \\ 0& I
    \end{bmatrix}
\end{align*}

We induct on the fact that $K - ww^*/{a_11}$ is itself an HPD matrix.  \textcolor{red}{Verify at your tables.}

Consequently, let's suppose we are at step $n$, and so far we have decomposed $A$ into the form 
$R_1*A_1R_1$.  If the upper left entry of the submatrix $K_1 - w_1w_1^*/a_{11}$, we can then write  $A_1 = \begin{bmatrix}1 & 0 \\ 0 & \tilde{R}_2^*\end{bmatrix}\begin{bmatrix}
        1 & 0 \\ 0 & \frac{1}{a_{11}}(K-ww^*)
    \end{bmatrix} \begin{bmatrix} 1 & 0 \\0 & \tilde{R}_2 \end{bmatrix}$,
    where $\tilde{R}_2$ is the matrix that allows us to apply the base case inductive step to $K- ww^*/a_{11}$.  Let $R_2 = \begin{bmatrix} 1 & 0 \\0 & \tilde{R}_2 \end{bmatrix}$.  Iterating, yields, eventually, 
    \begin{align*}
        A = R_1^*R_2^*& \cdots R_m^* R_m \cdots R_2R_1,
    \end{align*}
    as your book writes it.  It can also be written as 

\begin{theorem}
    Every hermitian positive definite matrix has a unique Cholesky factorization. 
\end{theorem}

Practice: Use QR, LU, SVD, and Cholesky to solve the least squares problem (see Lecture 12)

\end{document}



